%% ===========================================================================
%%  frontmatter/abstrak.tex — Abstrak
%%  Abstrak 200-300 kata. Di template ini abstrak memakai spasi 1; ganti
%%  \begin{singlespace} menjadi \begin{spacing}{2} bila pembimbing meminta
%%  spasi 2.
%% ===========================================================================
\babfrontmatter{Abstrak}

\begin{singlespace}
Penulisan karya ilmiah dengan bantuan \asing{large language model} (LLM)
menghadapi dua kelemahan yang berulang: sitasi yang dikarang dan penyuntingan
naskah yang berjalan tanpa sepengetahuan penulis. Laporan magang ini
menyajikan perancangan dan implementasi sistem penulisan riset berbantuan AI
yang dijalankan secara lokal pada komputer peneliti. Jawaban sistem hanya
bersumber pada dokumen yang dikumpulkan peneliti, sitasi hanya sah bila
terverifikasi terhadap pustaka lokal, dan draf tidak pernah masuk ke naskah
tanpa persetujuan manusia.

Sistem memadukan antarmuka web React, layanan \asing{backend} FastAPI dengan
sebelas perute API, basis data SQLite untuk data terstruktur, dan penyimpanan
vektor Chroma untuk potongan sumber setiap proyek. Dokumen sumber dicacah,
diubah menjadi vektor \asing{embedding}, dan disimpan, lalu dipakai pada tahap
pengambilan konteks. Percakapan berjalan dalam tiga mode: \asing{Ask} menjawab
pertanyaan
dari potongan sumber yang ditemukan dan menandai jawaban yang tidak berdasar,
\asing{Plan} hanya menyusun kerangka tanpa menulis naskah, dan \asing{Build}
menyusun draf bagian dari sumber. Konteks yang dikirim ke model dibatasi pada
potongan hasil pengambilan, ringkasan memori percakapan, dan konteks yang
disematkan peneliti. Alur revisi agen memetakan instruksi ke bagian naskah
melalui judul, alias, dan batas kata, lalu menggagalkan instruksi yang ambigu
alih-alih menebak. Setiap sitasi pada teks hasil revisi wajib terverifikasi
pada pustaka lokal yang bersumber dari OpenAlex, sehingga sitasi yang gagal
ditolak atau ditandai. Peneliti memutuskan menerima atau membuang setiap
usulan revisi. Naskah akhir diekspor sebagai berkas LaTeX beserta daftar
pustakanya dan disinkronkan dua arah dengan repositori GitHub untuk setiap
proyek. Pengujian dirancang dalam tiga skenario: keberpijakan sitasi,
ketepatan penargetan bagian naskah, dan ketepatan putaran ekspor serta
sinkronisasi.
% TODO: tambahkan angka hasil pengujian bila sudah tersedia.

\vspace{1\baselineskip}
\textbf{Kata kunci:} penulisan riset berbantuan AI, \asing{retrieval-augmented
generation}, verifikasi sitasi, \asing{human-in-the-loop}, OpenAlex,
\asing{large language model}
\end{singlespace}
